%% kapa: exam=1 round=33 year=2022 period=1 form=A subject=economics subject_ko=경제학원론 seq=30 number=70
%% source: Q-Net 감정평가사 기출문제 (cst00309) artlSeq=5209997 file=raw/1008/5209997/2022년도 제33회 감정평가사 1차시험 1교시(민법,경제학,부동산학).pdf
%% conversion: kordoc 4.24.51 -> tools/gen1.py (auto); review=proofread-sonnet; flags=none
\begin{question}{70}
자본이동이 완전한 소규모 개방경제의 먼델-플레밍(Mundell-Fleming)모형에서 변동환율제도인 경우, 긴축 통화정책을 시행할 때 나타나는 경제적 효과를 모두 고른 것은? (단, 물가수준은 고정이다.)
\begin{framedtext}
ㄱ. 소득 감소 \hfill ㄴ. 경상수지 개선 \hspace{6em} \par
ㄷ. 자국 통화가치 절하 \hfill ㄹ. 해외자본 유입 \hspace{6em}
\end{framedtext}
\begin{choices}
  \item ㄱ, ㄴ
  \item ㄱ, ㄷ
  \item ㄱ, ㄹ
  \item ㄴ, ㄷ
  \item ㄷ, ㄹ
\end{choices}
\end{question}
\answer{3}
